%document
\documentclass[10pt]{beamer}
%theme
\usetheme{metropolis}
% packages
\usepackage{color}
\usepackage{listings}
\usepackage[main=ngerman, USenglish]{babel}
\usepackage[utf8]{inputenc}
\usepackage{multicol}
\usepackage{csquotes}


% color definitions
\definecolor{mygreen}{rgb}{0,0.6,0}
\definecolor{mygray}{rgb}{0.5,0.5,0.5}
\definecolor{mymauve}{rgb}{0.58,0,0.82}
\definecolor{paleorange}{HTML}{CC7832}

\lstset{
    backgroundcolor=\color{white},
    % choose the background color;
    % you must add \usepackage{color} or \usepackage{xcolor}
    basicstyle=\footnotesize\ttfamily,
    % the size of the fonts that are used for the code
    breakatwhitespace=false,
    % sets if automatic breaks should only happen at whitespace
    breaklines=true,                 % sets automatic line breaking
    captionpos=b,                    % sets the caption-position to bottom
    commentstyle=\color{mygreen},    % comment style
    % deletekeywords={...},
    % if you want to delete keywords from the given language
    extendedchars=true,
    % lets you use non-ASCII characters;
    % for 8-bits encodings only, does not work with UTF-8
    frame=single,                    % adds a frame around the code
    keepspaces=true,
    % keeps spaces in text,
    % useful for keeping indentation of code
    % (possibly needs columns=flexible)
    keywordstyle=\color{blue},       % keyword style
    % morekeywords={*,...},
    % if you want to add more keywords to the set
    numbers=left,
    % where to put the line-numbers; possible values are (none, left, right)
    numbersep=5pt,
    % how far the line-numbers are from the code
    numberstyle=\tiny\color{mygray},
    % the style that is used for the line-numbers
    rulecolor=\color{black},
    % if not set, the frame-color may be changed on line-breaks
    % within not-black text (e.g. comments (green here))
    stepnumber=1,
    % the step between two line-numbers.
    % If it's 1, each line will be numbered
    stringstyle=\color{mymauve},     % string literal style
    tabsize=4,                       % sets default tabsize to 4 spaces
    % show the filename of files included with \lstinputlisting;
    % also try caption instead of title
    language = Java,
	showspaces = false,
	showtabs = false,
	showstringspaces = false,
	escapechar = ,
        morecomment=[s][\textcolor{paleorange}]{@}{\ },
}

\def\ContinueLineNumber{\lstset{firstnumber=last}}
\def\StartLineAt#1{\lstset{firstnumber=#1}}
\let\numberLineAt\StartLineAt



\newcommand{\codeline}[1]{
        \alert{\texttt{#1}}
}

\setmonofont{Fira Code}

% This Document contains the information about this course.

% Authors of the slides
\author{Marcus Köhler}

% Name of the Course
\institute{Java-Kurs}

% Fancy Logo
\titlegraphic{\hfill\includegraphics[height=1.25cm]{../templates/fsr_logo_cropped}}

\usepackage{csquotes}
\usepackage{tikz}
\usepackage{tikz-qtree}

\title{Java}
\subtitle{Scopes, Reference \& JavaDoc}
\date{\today}

\begin{document}

\begin{frame}
\titlepage
\end{frame}

\begin{frame}{Überblick}
    \setbeamertemplate{section in toc}[sections numbered]
    \tableofcontents
\end{frame}

\section{Scopes}

\begin{frame}[fragile]{Scopes}
    \center
    \large{Variablen sind nur innerhalb des Kontexts(d.h. innerhalb der \texttt{\{\}}) verfügbar, in denen sie deklariert wurden. \\
    \pause
    Das ist alles.}
    \pause
    \begin{columns}[T]
        \begin{column}{.5\textwidth}
            Richtig:
            \begin{lstlisting}[gobble=16]
                public class Scope {
                    int a = 5;
                    if(a <= 6) {
                        float f = (float)a;
                        System.out.println(f/1.25f);
                    }
                    System.out.println(a*10);
                }
            \end{lstlisting}
        \end{column}
        \begin{column}{.5\textwidth}
            Falsch:
            \begin{lstlisting}[gobble=16]
                public class Scope {
                    int a = 5;
                    if(a == 6) {
                        float f = (float)a;
                    }
                    System.out.println(f/1.25f);
                    //f is only available inside the if block
                    System.out.println(a*10);
                }
            \end{lstlisting}
        \end{column}
    \end{columns}
\end{frame}

\section{Die Reference}
\subsection{Intro}

\begin{frame}{Die Java-Reference}
    Die Java-Standardbibliothek ist in der sogenannten \textit{Reference} dokumentiert. Jede Java-Version hat ihre eigene Version der Reference;
    die momentan aktuellste Version(Java 10) ist unter \small\url{https://docs.oracle.com/javase/10/docs/api/overview-summary.html} \normalsize zu finden.\\
    \pause
    Die Reference ist hierarchisch gegliedert: \\
    \smallskip \small
    \begin{tikzpicture}
        \Tree   [.Reference
                    [.\node (level0) {\texttt{java.base}};
                        [.\node (level1) {\texttt{java.util}};
                            \node (level2) {\texttt{ArrayList}}; {\texttt{HashSet}} {\dots}
                        ]
                        [.{\texttt{java.lang}}
                            {\texttt{Object}} {\texttt{Math}} {\dots}
                        ]
                        {\dots}
                    ]
                    {\dots}
                ]
            \foreach \Value/\Text in {0/{Module},1/{Package},2/{Klasse}}
            {
              \node[anchor=west]
                at ([xshift=-2.5cm,yshift=0.05cm]{level2}|-{level\Value})
                {\Text};
            }
    \end{tikzpicture}
\end{frame}

\subsection{Mit der Reference arbeiten}
\begin{frame}{Struktur einer Reference-Seite}
    Die Seite einer einzelnen Klasse(oder eines Interfaces) in der Reference ist folgendermaßen strukturiert:
    \begin{enumerate}[<+->]
        \item Ein \enquote{Stammbaum}, d.h. wovon geerbt wird, wer erbt, was implementiert wird etc.
        \item Eine Beschreibung, wozu die Klasse/das Interface dient
        \item Ein Überblick über Attribute, Konstruktoren und Methoden, die \texttt{public} sind
        \item Detaillierte Beschreibungen von Konstruktoren und Methoden
    \end{enumerate}
\end{frame}

\begin{frame}{Die Java-Reference}
    Meistens ist das gesuchte in den detaillierten Methodenbeschreibungen zu finden. Darin stehen u.a.:
    \begin{itemize}[<+->]
        \item Die Parameter, die die Methode braucht
        \item Was eine Methode zurück gibt
        \item Welche Exceptions eine Methode generieren kann
    \end{itemize}
    \medskip
    \onslide<+->
    \begin{columns}
        \begin{column}{0.5\textwidth}
            Viele IDEs haben Möglichkeiten, die detaillierte Beschreibung einer Methode direkt anzuzeigen:\\
            IntelliJ: \texttt{Ctrl+Q} \\
            Eclipse: \texttt{F2} \\
            \smallskip
            Dabei muss der Cursor über dem Element sein, nachdem gesucht werden soll.
        \end{column}
        \begin{column}{0.5\textwidth}
            \includegraphics[width=\textwidth]{../images/IntelliJ_Reference_Popup.png}
        \end{column}
    \end{columns}
\end{frame}

\begin{frame}{Die Java-Reference}
    \center
    \huge
    Wäre doch auch praktisch, das für seinen eigenen Code zu haben, oder?
\end{frame}

\section{JavaDoc}
\subsection{Was ist JavaDoc?}
\begin{frame}[fragile]{Was ist JavaDoc?}
    JavaDoc(\textbf{Java Doc}umentation) ist ein Tool aus dem JDK, welches automatisch bestimmte Kommentare über
    Methoden, Konstruktoren und Attributen in das gleiche Format wie die Java-Reference umwandelt.\\
    \onslide<2->
    JavaDoc-Kommentare sind spezielle mehrzeilige Kommentare(\texttt{/* Kommentar */}), wobei der Anfang zwei Sterne (\texttt{/**}) am Anfang hat:
    \begin{lstlisting}[gobble=8]
        /**
         * Calculates the sum of two {@code ints}.
         *
         * @param a The first number of the equation.
         * @param b The second number of the equation.
         * @return A {@code long} that represents the sum of the two numbers.
         */
        public long sumOfInts(int a, int b) { ... }
    \end{lstlisting}
\end{frame}

\subsection{JavaDoc schreiben}
\begin{frame}{JavaDoc-Tags}
    JavaDoc hat ein paar spezielle Annotationen, die sogenannten \textit{Tags}. Diese dienen der Formatierung und beginnen immer mit \texttt{@}.\\
    \onslide<2->
    Einige oft verwendete Tags:\\
    \bigskip
    \begin{tabular}{l | l}
        \texttt{@author} & Gibt den Autor an \\
        \texttt{@param} \textit{<name>} & Beschreibung v. Param. \textit{<name>} (nur Methoden)\\
        \texttt{@return} & Beschreibung des Rückgabewertes (nur Methoden) \\
        \texttt{@throws} \textit{<exception>} & Wann \textit{<exception>} auftritt (nur Methoden)\\
        \texttt{\{@code} \textit{<text>}\texttt{\}} & Formatiert \textit{<text>} als Code(d.h. monospaced) \\
    \end{tabular} \\ %keep this linebreak, otherwise the bigskip underneath splits the following text in the middle
    \onslide<3->
    \bigskip
    Außer den JavaDoc-Tags kann man den Text auch mit einfachen HTML-Tags(wie <p>, <i> etc.) formatieren.
\end{frame}

\subsection{Das JavaDoc-Tool}
\begin{frame}[fragile]{Das JavaDoc-Tool}
    Um aus einer Quelldatei eine funktionierende Dokumentation zu generieren, muss das Programm \texttt{javadoc} aufgerufen werden.
    Ein Aufruf von \texttt{javadoc} kann folgendermaßen aussehen:
    \begin{lstlisting}[language=bash,gobble=8]
        # -d definiert das Zielverzeichnis
        javadoc org.test.example org.test.database -d ./doc/ #packages

        javadoc JavadocExample.java JavadocExercise.java -d ./doc/ #sourcefiles

        javadoc [package] [quelldateien] [optionen] #allgemein
    \end{lstlisting}
\end{frame}

\subsection{Übung}
\begin{frame}{Übung}
    \center
    \large
    Lade dir die Datei JavadocExercise.java aus dem Kurs-Repo herunter und vervollständige das Javadoc.
    Der genaue Inhalt sollte dabei die Funktion wiedergeben.
\end{frame}

\begin{frame}{Weitere Themen}
    \begin{itemize}[<+->]
        \item Eine Zusammenfassung und Wiederholung der Kurse seit der letzten Wiederholung
        \item Ein größeres praktisches Projekt(2 Termine)
        \item Ein kurzer Einblick hinter die Kulissen von Java(sehr theoretisch)
        \item Eine Einfühung zu Buildsystemen
        \item AuD-Inhalte in Java umgesetzt(z.B. Bäume, Quicksort, Dijkstra)
        \item Was sind Unit Tests und wie benutze ich sie?
        \item Vorschläge von euch
    \end{itemize}
\end{frame}
\end{document}
